% ===================================================================
%  اللوحة رقم ٩ — الجبل. الحرش.
%
%  Traduction, faite en 2026, en arabe levantin d'aujourd'hui, sur
%  l'ido de texto/io. Regles de la colonne : voir 10-lawha-01.tex.
%  Deux scenes, deux numerotations. Le plus gros tableau du livret :
%  47 blocs, 117 renvois. Un bloc marque « ¶2 » par ordo.py,
%  t09-c1-05-1 : une seule cle, deux alineas.
%
%  LE CEDRE (49) EST « الأرز », ET IL EST SUR LE DRAPEAU. C'est le
%  seul renvoi de tout le releve qui soit, dans la langue qui le
%  traduit, un embleme national : l'arbre du Mont-Liban, le bois du
%  temple de Salomon, celui que les rois d'Assyrie venaient couper et
%  dont il reste quelques centaines de pieds a Bcharre. La colonne
%  gujaratie avait « દેવદાર », le bois des dieux — le cedre de
%  l'Himalaya, autre espece, meme charge. DEUX COLONNES, DEUX CEDRES,
%  TOUS LES DEUX SACRES, et le livret les grave l'un et l'autre au
%  milieu d'une liste d'arbres ordinaires.
%
%  LES PLANTES DU LIVRET S'ARRETENT AU LEVANT, ET C'EST LA CINQUIEME
%  FOIS QUE CE TABLEAU-CI LE VERIFIE. Le sapin (11, 31) est
%  « الشوح », qui pousse sur ces montagnes ; le hetre (23) est
%  « الزان », le bouleau (26) « البتولا », le chene (29)
%  « السنديان », l'orme (47) « الدردار », la bruyere (57)
%  « الخلنج », la fougere (55) « السرخس », la pie (56) « العقعق ».
%  Aucune transcription. La colonne gujaratie avait du en transcrire
%  cinq dans ce meme tableau et la haoussa en decrire autant : LE
%  LIVRET A ETE DESSINE EN FRANCE, ET LE LEVANT EST L'AUTRE BOUT DE
%  LA MEME FLORE MEDITERRANEENNE.
%
%  MAIS LES BETES ALPINES, ELLES, NE SUIVENT PAS. La marmotte (39) est
%  « المرموط » et le chasseur de chamois (54) chasse « الشمواه » :
%  deux transcriptions, et ce sont les deux seules du tableau. LES
%  PLANTES S'ARRETENT AU LEVANT, LES ANIMAUX DE HAUTE MONTAGNE NON —
%  ce qui est exactement ce qu'on attend d'une flore qui longe une mer
%  et d'une faune qui suit une altitude.
%
%  LA TRUFFE A ETE EXAMINEE ET N'A PAS ETE REFUSEE, ET IL FAUT DIRE
%  POURQUOI. « الكمأة » est la truffe du desert, ramassee dans la
%  badiya syrienne apres les pluies de mars et vendue par sacs au
%  marche de Damas ; celle du livret est celle du Perigord, qu'un porc
%  deterre sous un chene. Le mot serait tombe sous la main comme
%  l'arachide gujaratie du tableau 7. Mais les refus de ce releve
%  separaient jusqu'ici des CATEGORIES — une legumineuse pour un
%  fruit a coque, un pavot a opium pour une mauvaise herbe de champ —
%  et ici il s'agit de deux genres du meme ordre, que le francais lui
%  meme appelle tous les deux truffes. LA REGLE A UN BORD, ET CE
%  RENVOI EST DESSUS : on ecrit « الكمأة ».
%
%  Deux mots pour finir. La lieue (alinea 15) est « فرسخ », le
%  parasang, vieille unite de marche que la geographie arabe emploie
%  depuis mille ans — le gujarati avait « ગાઉ » au meme endroit, et
%  ce sont deux fois la meme idee. Et la couleuvre (2) est
%  « الحنش », distincte de « الأفعى » la vipere, comme la limace (8)
%  « البزّاق » est distincte du « الحلزون » de l'escargot (4) : deux
%  paires que la colonne gujaratie n'avait pas pu faire.
% ===================================================================

%%K t09-apar-1 apar
\VUsaut{4.0mm}
\VUtitre{15.0pt}{60}{اللوحة رقم ٩}
\VUsaut{3.0mm}

%%K t09-tit-1 sub
\VUcentre{12.0pt}{0}{\textit{المشهد الأول.}}
\VUsaut{3.0mm}

%%K t09-tit-2 sub
\VUpk{11.4pt}{40}{الجبل. --- بلدة المياه المعدنية.}
\VUsaut{3.0mm}

%%K t09-c1-01-1 p
1. --- صار لي تمن أيام ببراتز. وما فيني وصّف كل الإعجاب اللي حسّيته
لمّا وقفت قدّام \VUgras{سلسلة الجبال} \textsuperscript{(1)} العجيبة تبع
البيريني، اللي \VUgras{خطّ قممها} \textsuperscript{(2)} بيبيّن عالسما
الزرقا متل زخرفة عملاقة.

%%K t09-c1-02-1 p
2. --- ما كنت متصوّر إنّه \VUgras{قمم} \textsuperscript{(3)} البيريني
بتوصل لهالـ\VUgras{علوّ} الكبير، وضلّيت مذهول قدّام \VUgras{الأنهار
الجليدية} \textsuperscript{(5)} الرائعة و\VUgras{رؤوس الجبال}
\textsuperscript{(4)} المتلّجة اللي عم تتدحرج عنها \VUgras{انهيارات
تلجية} مرعبة.

%%K t09-tit-3 sub
\VUsaut{3.0mm}
\VUpk{10.2pt}{40}{منظر الجبل.}
\VUsaut{3.0mm}

%%K t09-c1-03-1 p
3. --- من تاني يوم وصولي، رحت عالـ\VUgras{بركان} \textsuperscript{(6)}
العتيق الخامد اللي جنب براتز. وعلى حواف \VUgras{فوّهته} اليابسة
والعارية، لسّا منشوف منيح آثار مرقان \VUgras{الحمم} اللي جرت، من كم
قرن، على \VUgras{سفح الجبل} \textsuperscript{(7)}. ولقيت، على وحدة من
\VUgras{المنحدرات}، كم قطعة من هالحمم.

%%K t09-c1-04-1 p
4. --- وبراتز عالحدود، و\VUgras{القلعة} \textsuperscript{(8)} اللي بتحمي
\VUgras{المضيق} \textsuperscript{(27)} اللي بيفصل فرنسا عن إسبانيا،
بتذكّر تماماً بهديك القلاع العتيقة على ضفّة الراين اللي مبيّن إنّها عم
تترصّد \VUgras{فريسة} من فوق صخرها.

%%K t09-c1-05-1 p
5. --- \VUgras{نسر} \textsuperscript{(9)} ضخم عشّه مش بعيد عن
\VUgras{القلعة} \textsuperscript{(8)} كتير مرّات بيشدّ انتباهي.

هالـ\VUgras{طير الجارح}، اللي \VUgras{منقاره} قوي، كتير مرّات بيجيب
لصغاره خروف زغير أو جدي بيمسكه بـ\VUgras{مخالبه}. و\VUgras{جنحانه}
كبار لدرجة إنّه بيقدر يشحّط بالهوا وقت طويل مع حمله التقيل.

%%K t09-tit-4 sub
\VUsaut{3.0mm}
\VUpk{10.2pt}{40}{السايح.}
\VUsaut{3.0mm}

%%K t09-c1-06-1 p
6. --- وأحلى لفّة هونيك هيّي الرحلة على شلّال \textsuperscript{(10)}
« كاو ». \VUgras{الشلّال} بيوقع بدوّامات وبعدين بيختفي متل \VUgras{ساقية}
حلوة بـ\VUgras{حرش الشوح} \textsuperscript{(11)} اللي غربي المدينة.
وطلعنا من هالحرش لحدّ \VUgras{مرصد الأرصاد} \textsuperscript{(12)}، ومن
راس \VUgras{المنظرة}، شفنا \VUgras{المنظر الشامل} لكل المنطقة.
و\VUgras{المنظر} رائع، و\VUgras{الطبيعة} مبيّن إنّها عم تغدق عالبلد كل
الكنوز اللي عندها.

%%K t09-c1-07-1 p
7. --- وللنزول، أخدنا \VUgras{التلفريك} \textsuperscript{(13)} ووقفنا
بـ\VUgras{محطّة} زغيرة جنب \VUgras{خرايب} \textsuperscript{(14)}
\VUgras{قلعة إقطاعية} عتيقة كان ساكن فيها زمان أمراء براتز.

%%K t09-c1-08-1 p
8. --- مسكينة القلعة ! شو يمكن عم تفكّر وهيّي عم تشوف تحت
\VUgras{بلدة المياه المعدنية} \textsuperscript{(15)} الأنيقة اللي صارت
هلّق \VUgras{مصيف حمّامات} عالموضة ؟

%%K t09-c1-08-2 p
و\VUgras{المناخ} كتير حلو، و\VUgras{المياه المعدنية} كتير مفيدة، لدرجة
إنّه \VUgras{العلاج} اللي بينعمل بـ\VUgras{المصحّ} \textsuperscript{(17)}
متعة حقيقية.

%%K t09-tit-5 sub
\VUsaut{3.0mm}
\VUpk{10.2pt}{40}{بلدة حمّامات حلوة.}
\VUsaut{3.0mm}

%%K t09-c1-09-1 p
9. --- وبعدين، عملوا كل شي ليخلّوا الإقامة حلوة. \VUgras{جسر عالي}
\textsuperscript{(16)} كبير انبنى ليصير في سكّة حديد ؛
و\VUgras{منتزه} \textsuperscript{(18)} \VUgras{دار الحمّامات}، اللي
\VUgras{بيصير فيه حفلات}، مرتّب بشكل ساحر. وكم \VUgras{« فيلّا »}
\textsuperscript{(19)} رشيقة انبنوا حوالين \VUgras{الكازينو}
\textsuperscript{(21)}، و\VUgras{طريق معبّد} \textsuperscript{(22)} مصان
كتير منيح عم يسهّل المواصلات كتير.

%%K t09-c1-10-1 p
10. --- و\VUgras{جرف} \textsuperscript{(23)} بسيط بيفصل \VUgras{الطريق}
\textsuperscript{(20)} عن \VUgras{الوادي} \textsuperscript{(25)} نفسه، اللي
بقاعه \VUgras{بحيرة} \textsuperscript{(26)} زغيرة ورشيقة عم تغذّي
\VUgras{ساقية} \textsuperscript{(24)} صافية.

%%K t09-c1-11-1 p
11. --- وعالجهة التانية من الوادي، عم يشتغلوا \VUgras{منجم حديد}
\textsuperscript{(28)}. وكم مرّة رحت شفت \VUgras{عمّال المنجم} عم ينزلوا
بـ\VUgras{البير}، وحتى فتّ عالـ\VUgras{نفق} اللي بيقضّوا فيه نهارهن،
وهنّي عم يطلّعوا \VUgras{الخام} اللي فيه \VUgras{المعدن}.

%%K t09-c1-12-1 p
12. --- و\VUgras{مسبك} كبير انبنى جنب المنجم ليستعملوا حالاً الثروات
اللي عم تطلع منه. و\VUgras{الفرن العالي} \textsuperscript{(29)}، اللي عم
يتسخّن بـ\VUgras{الفحم الحجري} اللي كتير موجود هون، عم يخلّي الواحد
يحصّل \VUgras{الحديد المصبوب} بسرعة وبرخص.

%%K t09-c1-13-1 p
13. --- ومش بعيد عن المنجم عم يحفروا \VUgras{مقلع} \textsuperscript{(30)}.
و\VUgras{الحجّار} العتيق ما عم يقطع بـ\VUgras{كزمته} لا
\VUgras{غرانيت}، ولا \VUgras{سمّاق}، ولا حتى \VUgras{حجر رملي}، بس
\VUgras{حجر مقصوب} عادي لبنا البيوت.

%%K t09-c1-14-1 p
14. --- ومن أحلى زينة هالبلد \VUgras{الشوح} \textsuperscript{(31)}. بكل
محلّ بقاع \VUgras{وادي ضيّق} \textsuperscript{(32)}، على ضفّة
\VUgras{سيل} \textsuperscript{(33)}، أو \VUgras{هاوية}
\textsuperscript{(34)} أو جرف، \VUgras{إبره} الرفيعة عم تلوّن بالأخضر.

%%K t09-tit-6 sub
\VUsaut{3.0mm}
\VUpk{10.2pt}{40}{المشي عالأقدام.}
\VUsaut{3.0mm}

%%K t09-c1-15-1 p
15. --- عم استفيد من عطلتي لـ\VUgras{لفّ طويل}. و\VUgras{الدروب}
\textsuperscript{(35)} اللي عم تتلوّى عالجبل بتخلّي الواحد يقطع، بلا ما
ينتبه للمسافة، كم \VUgras{كيلومتر} وكتير مرّات كم \VUgras{فرسخ}.

%%K t09-c1-15-2 p
و\VUgras{الأنصاب} \textsuperscript{(36)} و\VUgras{أعمدة الإرشاد}
\textsuperscript{(37)} عم تعلّم الدروب اللي كتير مرّات بيلاقي عليها
الواحد شاب \VUgras{جبلي} \textsuperscript{(38)} و\VUgras{الخرج}
\textsuperscript{(40)} على جنبه، شايل \VUgras{مرموط}
\textsuperscript{(39)}، أو كم \VUgras{غجري} \textsuperscript{(41)}
\VUgras{طاقية فروه} \textsuperscript{(42)} بتتناقض بشكل غريب مع
\VUgras{قبّوعه} \textsuperscript{(43)} العتيق المشرشح. وهوّي عم يسوق
بسلسلة \VUgras{دبّ} \textsuperscript{(44)} مدرّب.

%%K t09-tit-7 sub
\VUpk{10.2pt}{40}{تسلّق الجبال.}
\VUsaut{3.0mm}

%%K t09-c1-16-1 p
16. --- كل يوم تقريباً بروح مع \VUgras{دليل} \textsuperscript{(48)}
لتمرّن عالـ\VUgras{تسلّق}، وبفتخر كتير لمّا، و\VUgras{الكيس}
\textsuperscript{(47)} على ضهري و« \VUgras{عصا التسلّق} » بإيدي، ومتل
\VUgras{سايح} \textsuperscript{(46)} حقيقي، بهاجم \VUgras{الصخور}
\textsuperscript{(45)}.

%%K t09-c1-17-1 p
17. --- ومن راس الجبل، سكّان السهل ما بيبيّنوا أكبر من النمل ؛
و\VUgras{قطيع الغنم} \textsuperscript{(50)} اللي عم يرعى بـ\VUgras{المرعى}
\textsuperscript{(49)} بيبيّن متل لعبة ولاد. و\VUgras{الصنوبرة}
\textsuperscript{(51)} أو \VUgras{بيت الخشب} \textsuperscript{(52)} بالكاد
بيبيّنوا متل نقطة عالخضرة.

%%K t09-c1-18-1 p
18. --- وبهالرحلات، كتير مرّات منلاقي عالـ\VUgras{درب}
\textsuperscript{(53)} \VUgras{صيّاد الشمواه} \textsuperscript{(54)} شايل
عكتفه الحيوان اللي هلّق قتله.

%%K t09-c1-19-1 p
19. --- وحطّيت بمجموعة نباتاتي غصن \VUgras{سرخس} \textsuperscript{(55)}
كتير لافت وغصين \VUgras{خلنج} \textsuperscript{(57)} وردي حلو، قطفتهن عند
مدخل \VUgras{مغارة} \textsuperscript{(56)} زغيرة بآخر لفّة إلي.

%%K t09-tit-8 sub
\VUsaut{3.0mm}
\VUcentre{12.0pt}{0}{\textit{المشهد التاني.}}
\VUsaut{3.0mm}

%%K t09-tit-9 sub
\VUpk{11.4pt}{40}{الحرش. --- الصيد.}
\VUsaut{3.0mm}

%%K t09-c2-01-1 p
1. --- مبارح قضّيت نهار لذيذ بالحرش. والنشاط اللي بين الحيوانات
و\VUgras{الحشرات} \textsuperscript{(5)} اللي ساكنة فيه شوّقني كتير.

%%K t09-c2-01-2 p
\VUgras{العنكبوت} \textsuperscript{(1)}، اللي عم ينسج \VUgras{بيته}
\textsuperscript{(2)} بين غصنين ليمسك \VUgras{الدبّانة} \textsuperscript{(3)}
اللي عم تمرق ؛ و\VUgras{الحلزون} \textsuperscript{(4)}، اللي عم يشيل
\VUgras{قوقعته} بمشقّة ؛ و\VUgras{الخنفسا أبو قرون}، اللي عم تفتّش على
فريستها ؛ والنملة النشيطة، اللي كتير حماسية جنب \VUgras{بيت النمل}
\textsuperscript{(6)} — كل هالمخلوقات الزغيرة كانت رايحة وجايي وعم
تشتغل بنشاط ما إلو متيل.

%%K t09-c2-02-1 p
2. --- و\VUgras{حيّة} \textsuperscript{(7)}، حسبتها بالنظرة الأولى
\VUgras{أفعى}، بس ما كانت غير \VUgras{حنش} عادي، كانت مختبّية تحت جذع
شجرة، ومن وقت لوقت عم تطلّع راسها الرفيع اللي دايماً مبيّن جاهز
ليعضّ. وقدّامي \VUgras{بزّاق} \textsuperscript{(8)} كبير عم يزحف بثقل بعد
ما أكل وحدة من \VUgras{الفريز البرّي} \textsuperscript{(11)} الطيّب اللي
عم ينبت هونيك بكترة.

%%K t09-c2-03-1 p
3. --- والأرض كانت مفروشة بـ\VUgras{زهور الحرش} ؛ \VUgras{زهرة الربيع}
\textsuperscript{(9)}، و\VUgras{الونكة} \textsuperscript{(10)}، وكتير نباتات
تانية ما كنت بعرفها، كانوا مزهّرين بين \VUgras{خصل العشب}
\textsuperscript{(12)}.

%%K t09-c2-04-1 p
4. --- وبهالمنطقة، \VUgras{الكمأة} تقريباً بكترة \VUgras{الفطر}
\textsuperscript{(14)}، و\VUgras{خنّوص} \textsuperscript{(13)} تبع حارس
الحرش كان عم يفتّش بشراهة إذا رح يلاقي منها شوي.

%%K t09-tit-10 sub
\VUsaut{3.0mm}
\VUpk{10.2pt}{40}{الصيد المخالف.}
\VUsaut{3.0mm}

%%K t09-c2-05-1 p
5. --- حضرت \VUgras{آخر} مشهد سلّاني كتير. وبمساعدة شمّ \VUgras{كلبه
القصير} \textsuperscript{(15)}، \VUgras{حارس الحرش} \textsuperscript{(16)}
هلّق فاجأ \VUgras{صيّاد مخالف} \textsuperscript{(22)}، بالضبط لمّا كان عم
يحضّر حاله يشيل \VUgras{الصيد} \textsuperscript{(17)} اللي قتله. ولأنه
فضّل يترك صيده على إنّه ينمسك، الصيّاد المخالف هرب، و\VUgras{الأرنب
البرّي} \textsuperscript{(18)}، و\VUgras{الغزالة} \textsuperscript{(19)}،
و\VUgras{الوعل الزغير} \textsuperscript{(20)} و\VUgras{الخنزير البرّي}
\textsuperscript{(21)} ضلّوا مرميين عالعشب وعم يفرّحوا حارس الحرش.

%%K t09-c2-06-1 p
6. --- وتنوّع الشجر اللي بالحرش شي بيدهش. \VUgras{الزان}
\textsuperscript{(23)} و\VUgras{البتولا} \textsuperscript{(26)}،
و\VUgras{الصنوبر} اللي بيعطي \VUgras{الراتينج} \textsuperscript{(27)}،
و\VUgras{السنديان} \textsuperscript{(29)} وكمان كتير غيرهن عم يخلطوا
أغصانهن.

%%K t09-c2-06-2 p
و\VUgras{اللبلاب} \textsuperscript{(24)}، اللي عم يلتصق بجذوع الشجر
العالي، عم يلوّن \VUgras{الغابة} \textsuperscript{(25)} بخضرة لامعة، عم
تتّحد بشكل حلو مع اللون الأخضر الأفتح والأبهت تبع \VUgras{الطحلب}
\textsuperscript{(31)} اللي عم يدوّر فيه \VUgras{نقّار الخشب}
\textsuperscript{(30)} عالحشرات.

%%K t09-tit-11 sub
\VUsaut{3.0mm}
\VUpk{10.2pt}{40}{منظر الحرش.}
\VUsaut{3.0mm}

%%K t09-c2-07-1 p
7. --- والسنديان العتيق سحرني. \VUgras{جذوره} \textsuperscript{(32)}
الغليظة والملتوية، الطالعة نصّها من الأرض، عم تعطيه منظر رائع من قوّة
وعنفوان، وعم سأل حالي كيف \VUgras{صاحب الملك} \textsuperscript{(35)}
بيقدر يرضى يقطعهن ويسلّمهن لـ\VUgras{المنشار} \textsuperscript{(34)} تبع
\VUgras{الحطّاب} \textsuperscript{(33)}. و\VUgras{الجذوع}
\textsuperscript{(36)} المساكين، المسلوخين من \VUgras{قشرهن}
\textsuperscript{(37)} أو المشقوقين بـ\VUgras{فاس} \textsuperscript{(38)}
\VUgras{العامل باليومية} \textsuperscript{(39)}، عم يحزّنوني.

%%K t09-c2-08-1 p
8. --- وجنبه، \VUgras{فحّام} \textsuperscript{(41)} عتيق حضّر
بـ\VUgras{مجرفته} \textsuperscript{(42)} كومة فحم، وختيارة شالت
\VUgras{حزمة} \textsuperscript{(40)} \VUgras{حطب يابس} من غصون الشجر
المقطوع.

%%K t09-tit-12 sub
\VUsaut{3.0mm}
\VUpk{10.2pt}{40}{الصيد.}
\VUsaut{3.0mm}

%%K t09-c2-09-1 p
9. --- كنت بدّي روح ارتاح كم لحظة بـ\VUgras{بيت حارس الحرش}
\textsuperscript{(46)}، بس لمّا وصلت جنب \VUgras{البركة}
\textsuperscript{(43)}، اللي عم تعوم عليها \VUgras{بجعة}
\textsuperscript{(45)} حلوة، لاحظت إنّه \VUgras{فيضان}
\textsuperscript{(44)} جديد حوّل الطريق لمستنقع حقيقي. واضطرّيت أغيّر
طريقي، ووأنا مارق جنب \VUgras{الأرز} \textsuperscript{(49)}،
و\VUgras{الحور} \textsuperscript{(48)}، و\VUgras{الدردار}
\textsuperscript{(47)} اللي على طرف الحرش، شفت طالع من \VUgras{أجمة}
\textsuperscript{(54)} \VUgras{أيّل} \textsuperscript{(50)} رائع، وراه
\VUgras{كلاب صيد} \textsuperscript{(51)} عنيدة اللي هيّجها كمان
\VUgras{بوق} \textsuperscript{(53)} \VUgras{الصيّادين عالخيل}
\textsuperscript{(52)}. والحيوان المسكين كان تعبان تماماً. وإجا وقع وهوّي
عم يموت على كم خطوة منّي.

%%K t09-c2-10-1 p
10. --- وابن حارس الحرش، اللي جذبته ضجّة \VUgras{المطاردة}، طلع من
البيت ورجعنا نحنا التنين عالحرش. كان شايف \VUgras{عقعق}
\textsuperscript{(56)} على غصن سنديانة عتيقة. وفكّر إنّه عشّه يمكن مخبّى
بـ\VUgras{الورق} \textsuperscript{(58)}، وبدّه يمسك العشّ.

%%K t09-c2-10-2 p
وخفيف متل \VUgras{السنجاب} \textsuperscript{(61)}، تسلّق من \VUgras{غصن}
\textsuperscript{(55)} لغصن، بس ما لقي شي وما قدر غير يطيّر
\VUgras{زاغ} \textsuperscript{(60)} عتيق كان حاطط على \VUgras{فرع}
\textsuperscript{(57)}.

\VUsaut{4.0mm}
\VUfilet{22mm}
